\section{Operational selection in the same Lorentz experiment}\label{sec:operational}
\subsection{An exact restriction with a nonzero realizability loss}
Use the whole equilibrium preparation and $T<d_*$. Let $S$ be the initial velocity direction. A direction-only preparation device may change the direction distribution, but keeps the position uniform on $Q_R$ conditional on that direction. This is an explicit restricted preparation class; it is not a claim about every possible apparatus.

\begin{proposition}[Direction-only preparation preserves the hit probability]\label{prop:direction-prep}
For every fixed initial direction $\vartheta$,
\[
 P_R(C_R=1\mid\vartheta)=p_R=2RT/A_R.
\]
Consequently any normalized direction-only reweighting leaves the hit probability unchanged. For the source $V=\xi C_R$, its restricted information-priced optimum is $\xi p_R$, attained by the original equilibrium law, and the exact deficit is
\begin{equation}
 \Delta_{\rm dir}(R,\xi)
 =\log(1-p_R+p_Re^\xi)-\xi p_R.
 \label{eq:direction-gap}
\end{equation}
It is strictly positive for $\xi\ne0$ and has expansion
$\frac12\xi^2p_R(1-p_R)+O(\xi^3)$.
\end{proposition}
\begin{proof}
At fixed $\vartheta$, use incoming boundary position and elapsed time as coordinates. The area element is $R|n_\alpha\cdot v_\vartheta|\dd\alpha\dd s$ on the incoming semicircle. Its integral over $0<s<T$ is
$RT\int_{-\pi/2}^{\pi/2}\cos\phi\dd\phi=2RT$.
No overlap occurs because there is at most one collision. Divide by $A_R$ to obtain the conditional probability, independent of direction. Any direction density therefore gives the same mean of $C_R$. For $V=\xi C_R$, the conditional mean given direction is the constant $\xi p_R$. Theorem~\ref{thm:coarse-prep} gives value $\xi p_R$ and optimizer $P_R$. The unrestricted exponential moment is $1-p_R+p_Re^\xi$, proving \eqref{eq:direction-gap}. Strict Jensen gives positivity; direct Taylor expansion gives the stated tangent.
\end{proof}

Thus a formally unique conjugate source does not by itself make a target hit fraction available to a direction-only preparation device. This is an exact calculation of a restricted class, not a claim that the desired selection is impossible with a richer protocol. Completed-record selection realizes the tilted family as follows.

\subsection{A parameter-independent postselection protocol}
For a prescribed $\xi\in\R$, observe the deterministic hit bit and accept the completed trial with probability
\[
 a_\xi(C_R)=\exp\{\xi C_R-\max(0,\xi)\}.
\]
Only the observed bit and an independently prepared threshold coordinate are needed; the acceptance formula does not use the unknown radius.

\begin{theorem}[Budgeted completed-record realization]\label{thm:lorentz-selection}
Under independently repeated equilibrium preparations and threshold coordinates, the accepted completed-record law is $Q_R^\xi\propto e^{\xi C_R}P_R$ for every $R\in I$. Its success probability is
\[
 s_R(\xi)=e^{-\max(0,\xi)}(1-p_R+p_Re^\xi)\ge e^{-|\xi|}.
\]
Using at most $B$ trials, conditional on obtaining an accepted record the output law is exactly $Q_R^\xi$, and the failure probability is at most
$\exp(-Be^{-|\xi|})$. The protocol realizes postselected records, not a changed real-time force on the original trial.
\end{theorem}
\begin{proof}
Proposition~\ref{prop:acceptance} gives the accepted law and success probability. Since $\xi C_R$ has oscillation $|\xi|$, its minimum acceptance is $e^{-|\xi|}$. For first acceptance at trial $j\le B$, independence gives joint probability $(1-s_R)^{j-1}s_RQ_R^\xi(A)$ of acceptance and an output in $A$. Sum in $j$ and divide by $1-(1-s_R)^B$ to get the exact conditional output law. The failure bound follows from $(1-s_R)^B\le e^{-Bs_R}$. The bit is measured only after the trajectory, so the construction is causal postselection and changes no earlier force.
\end{proof}

The trial count is a genuine operational budget. A bound on total mechanical work or reset cost additionally requires a cost bound for one preparation/readout/reset cycle. Under a separately verified deterministic cycle cost $W_*$, $B$ trials cost at most $BW_*$. Relative entropy does not establish $W_*$. Nor does one retained trajectory provide the independent re-preparations used here.

\subsection{Conjugate targets and calibration error}
The tilted hit probability is
\begin{equation}
 q_R(\xi)=\frac{e^\xi p_R}{1-p_R+e^\xi p_R}.
 \label{eq:qR}
\end{equation}
For a prescribed $c\in(0,1)$, the formal target source is
\[
 \xi_R(c)=\log\frac c{1-c}-\log\frac{p_R}{1-p_R},\quad
 \partial_R\xi_R(c)=-\frac{p_R'}{p_R(1-p_R)}.
\]
The minimum unconstrained information cost is
$c\log(c/p_R)+(1-c)\log((1-c)/(1-p_R))$.
These formulas follow by the Bernoulli entropy calculation. In a direction-only class the target is unavailable unless $c=p_R$; in completed-record postselection a fixed source is implementable without knowing $R$, but selecting the source for a specific target $c$ uses knowledge or an estimate of $p_R$.

\begin{proposition}[Source and estimated-probability error]\label{prop:calibration-error}
For the same reference mechanical experiment,
\[
 d_{\TV}(Q_R^\xi,Q_R^{\widetilde\xi})
 =|q_R(\xi)-q_R(\widetilde\xi)|\le\tfrac14|\xi-\widetilde\xi|.
\]
If $p_R,\widehat p\in[\epsilon,1-\epsilon]$ and
$\widehat\xi=\operatorname{logit}(c)-\operatorname{logit}(\widehat p)$, then
\[
 |q_R(\widehat\xi)-c|
 \le\frac{|\widehat p-p_R|}{4\epsilon(1-\epsilon)}.
\]
\end{proposition}
\begin{proof}
Tilting a binary variable preserves the two conditional path laws given that variable. Because those laws have disjoint supports, the TV distance of the full selected laws equals the TV distance of their Bernoulli weights. Differentiating \eqref{eq:qR} gives $\partial_\xi q_R=q_R(1-q_R)\le1/4$. For the second assertion, the logit function has derivative $1/[p(1-p)]\le1/[\epsilon(1-\epsilon)]$ on the interval; apply the first bound and use $q_R(\xi_R(c))=c$.
\end{proof}

For capped preparation, Theorem~\ref{thm:cap} gives the exact optimal density for any declared initial statistic. Its evaluation may depend on the candidate radius through $E_R[V\mid S]$. A radius-dependent mask is not automatically a valid unknown-parameter policy. Either the policy is fixed independently of $R$, or its calibration and error must be accounted for. The parameter-independent completed-bit rule above is an explicit way to avoid this hidden use of the true radius when the source, rather than the target mean, is prescribed.

\subsection{A finite information budget is not a unique physical law}
For bounded $V$, maximizing $QV$ subject to $D(Q\Vert P)\le b$ gives the bound
\[
 QV\le\frac{\log Pe^{tV}+b}{t},\qquad t>0,
\]
by Theorem~\ref{thm:gibbs}. Equality holds when $Q=Q^{tV}$ and its entropy equals $b$. The function $b(t)=t\Psi'(t)-\Psi(t)$ has derivative $t\Var_{Q^{tV}}V\ge0$. Thus an interior nondegenerate branch can be solved by a scalar entropy equation. Boundary budgets or unavailable preparation operations can prevent equality. This information constraint remains separate from a mechanical energy constraint, even in the same experiment.
\par
